\documentclass[12pt, a4paper]{article}

% --- Pacotes Básicos ---
\usepackage[utf8]{inputenc}
\usepackage[T1]{fontenc}
\usepackage[main=brazilian, english]{babel}
\usepackage{indentfirst}
\usepackage{geometry}
\geometry{left=3cm, top=3cm, right=2cm, bottom=2cm}

% --- Informações do Documento ---
\title{Resenha: Obsolescência Programada}
\author{Erickson Giesel Müller}
\date{\today}

\begin{document}
\maketitle
\thispagestyle{empty} % primeira página não tem número.

	O documentário "Obsolescência Programada" busca se aprofundar no conceito de estarmos sujeitos ao prazo de vencimento de produtos que são planejados para durarem apenas até um limite de tempo, com o intuito de manter a população consumindo num ciclo repetitivo. O filme busca abordar o conceito desde o seu surgimento, para isso, inicia com a história de uma lâmpada centenária que está instalada no edifício do corpo de bombeiros de uma cidadezinha dos Estados Unidos. Essa lâmpada funciona até hoje, mas ela foi fabricada logo antes de uma época que passaram a controlar a fabricação para que as lâmpadas durassem até 10 mil horas.

	Outra história que é abordada durante o documentário é o caso de um jovem cuja impressora doméstica parou de funcionar repentinamente. Ele vai atrás de diversos técnicos especialistas e todos eles recomendam que ele compre uma nova, pois o serviço de identificar o problema e consertá-lo seria mais caro que a própria impressora. Inconformado por essa situação, o jovem busca se aprofundar no problema e vai à internet para saber o que pode ter causado o erro repentino. Essa história é narrada em partes ao longo do filme, mas no final o homem descobre que existe um chip na placa da impressora que conta quantas impressões ela faz e interrompe o funcionamento após um certo número de impressões, o homem então baixa um \textit{software} que reinicia essa contagem e a impressora volta a funcionar normalmente.

	A obsolescência programada foi uma das soluções pensadas para a crise de 1929, diminuindo o prazo de vencimento de produtos, e impondo leis que desincentivem a população a possuir produtos após um período desde a fabricação, com a desculpa de ser por motivos de segurança, poderiam reascender o consumo da população e movimentar a economia. Essa solução não foi aplicada nessa época, mas aconteceu algo similar no modo de vida americano, os produtos passaram a adotar um design que desperte o desejo de consumo, produzidos num ciclo que faz o cidadão americano sentir desejo de comprar um veículo novo a cada três anos apenas por ser o "modelo do ano".

	Pode-se perceber que o modelo de produção dos produtos de hoje em dia não pensados tendo como centro o consumidor, mas sim o lucro do fabricante, poderíamos hoje ter produtos de qualidade que não demandam substituição frequente, mas isso não seria atrativo ao mercado. O resultado da obsolescência programada é um dano ambiental grave, todos os dias, toneladas de lixo são despejadas no continente africano. Desde a época em que estávamos fabricando lâmpadas que duram 100 anos até hoje, podemos perceber que o nível de consumismo da humanidade só vem aumentando, um dos maiores problemas é o dano ambiental e podemos identificar a obsolescência programada como um dos agentes causadores desse problema.

\end{document}
